\documentclass[letterpaper,11pt]{article}
\usepackage[T1]{fontenc}
\usepackage{latexsym}
\usepackage[empty]{fullpage}
\usepackage{titlesec}
\usepackage{marvosym}
\usepackage[usenames,dvipsnames]{color}
\usepackage{verbatim}
\usepackage{enumitem}
\usepackage[hidelinks]{hyperref}
\usepackage{fancyhdr}
\usepackage[italian]{babel}
\usepackage{tabularx}
\usepackage{mathptmx}

% Harvard-style layout from the supplied main.tex.
\pagestyle{fancy}
\fancyhf{}
\fancyfoot{}
\renewcommand{\headrulewidth}{0pt}
\renewcommand{\footrulewidth}{0pt}
\setlength{\footskip}{12pt}
\addtolength{\oddsidemargin}{-0.5in}
\addtolength{\evensidemargin}{-0.5in}
\addtolength{\textwidth}{1in}
\addtolength{\topmargin}{-.5in}
\addtolength{\textheight}{1.0in}
\urlstyle{same}
\raggedbottom
\raggedright
\setlength{\tabcolsep}{0in}
\titleformat{\section}{\vspace{-4pt}\scshape\raggedright\large}{}{0em}{}[\color{black}\titlerule\vspace{-5pt}]
\newcommand{\resumeItem}[1]{\item\small{{#1\vspace{-2pt}}}}
\newcommand{\resumeSubheading}[4]{
  \vspace{-1pt}\item
  \begin{tabular*}{0.97\textwidth}[t]{l@{\extracolsep{\fill}}r}
    \textbf{#1} & #2 \\
    \textit{\small#3} & \textit{\small#4} \\
  \end{tabular*}\vspace{-3pt}
}
\newcommand{\resumeProjectHeading}[2]{
  \item
  \begin{tabular*}{0.97\textwidth}{l@{\extracolsep{\fill}}r}
    \small#1 & #2 \\
  \end{tabular*}\vspace{-3pt}
}
\renewcommand\labelitemii{$\vcenter{\hbox{\tiny$\bullet$}}$}
\newcommand{\resumeSubHeadingListStart}{\begin{itemize}[leftmargin=0.15in,label={},topsep=3pt,itemsep=4pt,parsep=0pt]}
\newcommand{\resumeSubHeadingListEnd}{\end{itemize}}
\newcommand{\resumeItemListStart}{\begin{itemize}[topsep=3pt,itemsep=2pt,parsep=0pt]}
\newcommand{\resumeItemListEnd}{\end{itemize}\vspace{-5pt}}
\hypersetup{pdftitle={Davide Di Matteo - Curriculum Vitae (IT)},pdfauthor={Davide Di Matteo},pdflang={it-IT}}

\begin{document}
\begin{center}
  \textbf{\Huge\scshape Davide Di Matteo} \\ \vspace{1pt}
  Via Taranto 178, Roma, 00182 \\ \vspace{1pt}
  \small (+39) 346 134 5120 $|$
  \href{mailto:davidedimatteo97@gmail.com}{\underline{davidedimatteo97@gmail.com}} $|$
  \href{https://linkedin.com/in/davide-di-matteo}{\underline{linkedin.com/in/davide-di-matteo}}
\end{center}

\section{Formazione}
\resumeSubHeadingListStart
  \resumeSubheading
    {Università degli Studi di Roma ``Tor Vergata''}{Roma, Italia}
    {Laurea Triennale in Informatica}{Set 2016 -- Ott 2022}
    \resumeItemListStart
      \resumeItem{\textbf{Tesi:} ``Analisi di sicurezza e capacità dei software Anti Malware su Linux'' -- Sviluppo di tecniche di evasione malware (Process Injection, Memory Shellcode) in C e Python.}
    \resumeItemListEnd
  \resumeSubheading{Certificazioni Professionali}{}{Cybersecurity \& Networking}{}
    \resumeItemListStart
      \resumeItem{\textbf{Certified Ethical Hacker (CEH)} -- EC-Council (Nov 2025)}
      \resumeItem{\textbf{CompTIA Security+} -- CompTIA (Ott 2017)}
    \resumeItemListEnd
\resumeSubHeadingListEnd

\section{Esperienza Professionale}
\resumeSubHeadingListStart
  \resumeSubheading{Aeroporti di Roma (ADR)}{Roma, Italia}
    {Cyber Security Specialist}{Mar 2026 -- Presente}
    \resumeItemListStart
      \resumeItem{Coordinamento di un team dedicato al ciclo di Vulnerability Management dei sistemi aeroportuali.}
      \resumeItem{Gestione delle attività VAPT nel Perimetro di Sicurezza Nazionale Cibernetica (PSNC): coordinamento dei fornitori, monitoraggio delle remediation e retest.}
      \resumeItem{Supporto all'Incident Response con il CSOC e collaborazione con i team di Cyber Threat Intelligence.}
    \resumeItemListEnd
  \resumeSubheading{Leonardo S.p.A.}{Roma, Italia}
    {Cyber Security Assurance and Red Team}{Nov 2024 -- Feb 2026}
    \resumeItemListStart
      \resumeItem{Vulnerability Assessment e Penetration Testing su sistemi INFOSEC e applicazioni critiche con Nessus, Burp Suite e Nmap.}
      \resumeItem{Supporto alle certificazioni Common Criteria (ISO/IEC 15408), con documentazione tecnica e piani di remediation.}
      \resumeItem{Collaborazione con il Laboratorio Valutazione Sicurezza (LVS) nell'analisi della resistenza agli attacchi e nello sviluppo di procedure di test per il Laboratorio Accreditato di Prova (LAP).}
    \resumeItemListEnd
  \resumeSubheading{Sipal S.p.A.}{Roma, Italia}
    {Cybersecurity Analyst $|$ Operatore Ce.Va.}{Dic 2022 -- Nov 2024}
    \resumeItemListStart
      \resumeItem{Analisi dei requisiti di sicurezza e penetration test su prodotti dei clienti per le valutazioni Common Criteria.}
      \resumeItem{Analisi statica e dinamica (SAST/DAST) e vulnerability scanning per individuare vulnerabilità nel codice e nell'infrastruttura.}
      \resumeItem{Collaborazione con gli sviluppatori su mitigazione delle vulnerabilità e documentazione tecnica di valutazione.}
    \resumeItemListEnd
\resumeSubHeadingListEnd

\section{Progetti Tecnici e Ricerca}
\resumeSubHeadingListStart
  \resumeProjectHeading{\textbf{\href{https://dingo97.github.io/it/research/CVE-2026-20516/}{CVE-2026-20516 -- MediaTek}} $|$ \emph{Android TV}}{Set 2026}
    \resumeItemListStart
      \resumeItem{Ricercatore accreditato per una vulnerabilità \emph{confused deputy} in MiracastService (CVSS 5.5). Riproduzione via ADB, analisi del confine di fiducia e divulgazione coordinata con il vendor.}
    \resumeItemListEnd
  \resumeProjectHeading{\textbf{Linux Malware Evasion Research} $|$ \emph{C, Python}}{Tesi Sperimentale}
    \resumeItemListStart
      \resumeItem{Sviluppo di malware custom in C e Python per testare l'efficacia di sette antivirus Linux.}
      \resumeItem{Implementazione di Process Injection, Self-Injection e Direct Memory Shellcode Writing per valutare l'evasione dei controlli signature-based.}
    \resumeItemListEnd
  \resumeProjectHeading{\textbf{Piattaforma E-commerce Sicura} $|$ \emph{PHP, MySQL, HTML/CSS}}{Progetto Accademico}
    \resumeItemListStart
      \resumeItem{Sviluppo di un portale e-commerce full-stack con misure di sicurezza difensiva e gestione degli ordini; validazione dell'input per prevenire SQL Injection e XSS.}
    \resumeItemListEnd
  \resumeProjectHeading{\textbf{Database Management System (SoundCloud Clone)} $|$ \emph{MySQL, MongoDB}}{Progetto Accademico}
    \resumeItemListStart
      \resumeItem{Progettazione dello schema ER e implementazione di database relazionali e NoSQL ottimizzati per le prestazioni.}
    \resumeItemListEnd
\resumeSubHeadingListEnd

\section{Competenze Tecniche}
\begin{itemize}[leftmargin=0.15in,label={},topsep=3pt]
  \small\item{
    \textbf{Vulnerability Assessment \& Pentesting}: Nessus, Burp Suite, Nmap, Metasploit, SQLmap, Wireshark, OWASP Top 10 \\
    \textbf{Standard di Sicurezza}: Common Criteria (ISO/IEC 15408), metodologia CEH, PTES \\
    \textbf{Programmazione}: C, Python (Scripting \& Exploit Dev), Bash, PHP, SQL, HTML/CSS, JavaScript \\
    \textbf{Sistemi Operativi}: Linux (Admin \& Hardening, Kali, Ubuntu), Windows \\
    \textbf{Lingue}: Italiano (Madrelingua), Inglese (B2 - Tecnico)
  }
\end{itemize}
\end{document}
